% Root proof promoted after full monodromy and obstruction/crystal independent audits.
\subsection{A nonproduct three-neighbour and its clean theta boundary}

\begin{lemma}[A nonproduct principally polarized three-neighbour]
\label{lem:root-nonordinary-nonproduct-three-neighbor-candidate}
\status{Proved; both complete independent mathematical audits passed.}
Let $(A,\Theta_A)$ be a nonproduct nonordinary principally polarized
abelian surface over an algebraically closed field of characteristic two.
There is a nonproduct principally polarized surface $(B,\Theta_B)$
and a polarized isogeny $\phi:B\to A$ such that
\[
 \deg\phi=9,\qquad \ker\phi\subset B[3],\qquad
 \phi^*\lambda_{\Theta_A}=3\lambda_{\Theta_B}.
\]
The surface $B$ has the same two-rank as $A$. In two-rank zero,
every quotient of $A$ by a maximal isotropic subgroup of $A[3]$
is nonproduct. In two-rank one, at most one such subgroup has a
product principally polarized quotient.
\end{lemma}
\begin{proof}
Use the principal alternating Weil pairing on the etale group
$A[3]\simeq\mathbf F_3^4$. A maximal isotropic subgroup $H$ has
order nine. The polarized quotient construction already proved in
the odd-isogeny argument gives
\[
 \lambda:A\longrightarrow B=A/H,\qquad
 \lambda^*\lambda_{\Theta_B}=3\lambda_{\Theta_A}.
\]
Its polarized dual $\phi:B\to A$ satisfies
$\phi\lambda=[3]_A$, $\lambda\phi=[3]_B$ and the displayed
pullback formula. Both maps are etale. In particular they induce
isomorphisms on the two-divisible groups and preserve two-rank.

Suppose first that this rank is zero. On a nonproduct principal
surface the smooth theta curve has genus two. Its previously
proved characteristic-two normal form has equation $y^2+y=f_5(x)$;
the regular forms $dx,x\,dx$ have Cartier images $0,dx$.
Thus Cartier on its two-dimensional differential space has rank one.
The restriction of invariant differentials from its Jacobian is
an isomorphism. A rank-zero product principal surface is a product
of two supersingular elliptic curves, and its Cartier operator is
zero. Pullback by the etale isogeny $\lambda$ is an isomorphism on
invariant differentials and commutes with Cartier. Its target cannot
therefore be such a product.

Now suppose the rank is one and a quotient $B$ is a product
$E_o'\times E_s'$ with its product principal polarization. Exactly
one elliptic factor is ordinary and the other supersingular.
Let $E_o,E_s\subset A$ be the connected inverse-image elliptic
subvarieties of these axes. The restrictions of $\lambda$ are
isogenies onto the axes, so their respective types are ordinary
and supersingular. There is at most one elliptic subvariety of
each type in $A$: two distinct elliptic subvarieties span $A$ and
their addition map is an isogeny; two ordinary ones would give
two-rank two and two supersingular ones two-rank zero.

For either $E=E_o,E_s$, the projection formula for the restriction
of the polarized quotient gives
\[
             \deg(\lambda|_E)=3(\Theta_A\cdot E).
\]
Its kernel is a subgroup of $E[3]$, so this degree is one, three,
or nine. The intersection $\Theta_A\cdot E$ is a positive integer.
It cannot be one on a nonproduct principal surface. Indeed, in that
case the restricted polarization $\lambda_E:E\to E^\vee$ is an
isomorphism, and for the inclusion $i:E\hookrightarrow A$ the map
\[
 r=\lambda_E^{-1}\circ i^\vee\circ\lambda_{\Theta_A}:A\longrightarrow E
\]
satisfies $ri=1_E$. The maps $(e,z)\mapsto i(e)+z$ and
$a\mapsto(r(a),a-i(r(a)))$ are inverse group-scheme isomorphisms
$E\times\ker r\simeq A$. Hence $\ker r$ is a connected smooth
elliptic curve. The kernel is orthogonal to $E$ for the principal
polarization, which therefore becomes the product of the two
principal polarizations. This contradicts the nonproduct hypothesis.

It follows that $\Theta_A\cdot E=3$ and
$\deg(\lambda|_E)=9$. Thus $H=E[3]$ for either axis.
The uniqueness of the ordinary and supersingular elliptic
subvarieties shows that this possible bad subgroup $H$ is unique.
There are at least two maximal isotropic subgroups: a symplectic
basis supplies the spans of its two $e$ vectors and its two $f$
vectors. Choose one different from the possible bad subgroup.
Its quotient and polarized dual give the required nonproduct $B$.
\end{proof}

\begin{lemma}[A clean common three-theta boundary]
\label{lem:root-nonordinary-three-theta-clean-boundary-candidate}
\status{Proved; both complete independent mathematical audits passed.}
Let $(B,\Theta)$ be a nonproduct nonordinary principally polarized
abelian surface in characteristic two, and let $K\subset B[3]$
have order nine. Put $q:B\to B/K$. There are theta translates
$C_i=\Theta+a_i$, $1\le i\le3$, with $a_1+a_2+a_3=0$, such that:
each $q(C_i)$ is integral with eight ordinary nodes; the three
images meet pairwise transversely; and all self-nodes and mutual
intersections are distinct. For every $r_i\in K$ with
$r_1+r_2+r_3=0$, the source union
\[
                       C_r=\bigcup_{i=1}^3(C_i+r_i)
\]
belongs to the same complete system $|3\Theta|$, has six ordinary
intersection nodes, and has the same target image $Y$.
All these finitely many source unions lie in a connected coefficient
open in $|3\Theta|$ on which the additional identifications under
$q$ form finite etale node schemes, with eighteen nodes per
unordered nonzero type. Thus the first geometric hypothesis of
Theorem~\ref{thm:three-theta-conditional-full-image-candidate}
holds for this kernel.
\end{lemma}
\begin{proof}
Individual theta norms have the asserted eight nodes by
Corollary~\ref{cor:nonordinary-theta-three-torsion-nodal-norm-candidate}.
This also excludes triple identifications within one theta curve.
We choose the translates in the irreducible four-dimensional
parameter space $B^2$, with $a_3=-a_1-a_2$.

The locus of differences giving a tangency between two distinct
theta translates is contained in the image of the one-dimensional
theta curve under multiplication by two, together with the zero
difference. To see this, identify the theta curve with its
Abel--Jacobi image based at a Weierstrass point. Equality of tangent
lines is equality of canonical images, so the two points coincide
or are hyperelliptic conjugates. Their differences are respectively
zero or $2(P-W)$. Each pair-difference map from $B^2$ to $B$ is
surjective, including pairs involving $a_3$; one of its two
coefficients is $1$ or $-1$. Avoiding the finitely many translates
of this proper locus by $K$ makes all cross intersections transverse
and keeps distinct components distinct.

For three distinct components, a triple intersection after any
specified kernel translations is parametrized by the inverse image
of $\Theta^3$ under the affine translate of the homomorphism
\[
 B^3\longrightarrow B^3,\qquad
 (P,a_1,a_2)\longmapsto(P-a_1,P-a_2,P+a_1+a_2).
\]
Its integral coefficient matrix has determinant three, so it is
a finite etale isogeny in characteristic two. This incidence has
dimension three, and its image in the four-dimensional parameter
space is proper, since the projection is proper. There are only
finitely many kernel choices. For a triple involving two branches
from one component and one from another, the first two branches
come from one of finitely many fixed self-identifications on the
theta curve. Requiring a translate of another theta curve to contain
that point is again a proper condition: its pair difference lies
on a fixed translate of a curve. These exclusions also keep every
self-node off the other components and prevent coincidences among
the different mutual intersections. A finite union of proper closed
subsets cannot cover $B^2$. Choose a point in the remaining open.

The theorem of the square and $\sum a_i=\sum r_i=0$ give the actual
line-bundle identity $\mathcal O_B(C_r)\simeq\mathcal O_B(3\Theta)$.
The selected three pairs of theta translates each meet transversely
in two points, and no triples remain, so $C_r$ has exactly six nodes.
Translations in $K$ do not change its image under $q$.

Here is the required family statement near these reducible sources.
For fixed $g\in K\setminus\{0\}$, the two equations defining a
point of $C_s\cap(C_s+g)$ have independent differentials at every
point of each chosen $C_r$: both source branches are smooth there,
and their translated tangent lines are transverse. A source
intersection node cannot occur in such a pair, since it would give
at least three target branches, already excluded. The incidence
over the coefficient space is therefore etale at these points;
properness and quasi-finiteness make it finite etale after shrinking
to an open containing the finitely many chosen sources. Each fibre
has eighteen points, since $(3\Theta)^2=18$.

Impose the same open conditions for all nonzero $g$, exclude triple
identifications and component stabilizers, and identify $g$ with
$-g$. A node of unordered type $\{g,-g\}$ is represented exactly
once in the incidence for one chosen representative $g$; exponent
three and absence of triples ensure this assertion. The resulting
node scheme for each type has degree eighteen. The intersection of
these opens contains every chosen $C_r$ and is a nonempty open of
the irreducible coefficient space, hence connected. The usual
partial-normalization exact sequence applies there: the finite map
$C_s\to q(C_s)$ is a local isomorphism away from these pairs, and
at each pair it identifies two smooth transverse branches. This
gives the finite etale graph subgroups even when $C_s$ itself is
nodal. It proves precisely the stated geometric hypothesis.
\end{proof}

\begin{remark}[The remaining monodromy input]
\label{rem:root-three-neighbor-boundary-scope-candidate}
\status{Proved scope statement.}
These two lemmas supply a nonproduct source and a clean common
boundary for every fixed nonproduct nonordinary principal target
in degree three. They do not supply the product node-permutation
group or the gluing translations. Those are the second hypothesis
of the conditional graph theorem and require a separate proof.
No existence statement for clean theta norms at higher odd kernel
exponent is inferred.
\end{remark}
